% GENERATED FILE. Edit canonical lessons, then run npm run generate:books.
% canonical-source-hash: fnv1a64:3846f53f64307d13
% canonical-lessons: LA-C456-inquilinus, LA-C456-oppidanus, LA-C456-incola, LA-C456-incolo, LA-C456-habitator

\chapter{Tenant, Townsman, Resident, To live in, Occupant}
\label{ch:la-tenant-townsman-resident-to-live-in-occupant}
\hlchaptermodality{456}

\begin{chapteropening}
\textbf{By the end of this chapter:} \emph{I can talk about the people who live in a town: a tenant, a townsman, a resident, living in a place, and the occupant of a house.}

Complete the last lesson of chapter 456: I can talk about the people who live in a town: a tenant, a townsman, a resident, living in a place, and the occupant of a house.
\end{chapteropening}

\section[inquilīnus, inquilīnī]{inquilīnus, inquilīnī — a tenant, a lodger}

\label{lesson:LA-C456-inquilinus}



\begin{quote}
Before the new one: say the Latin for to delay, then the Latin for to stop for a while.
\end{quote}

\subsection*{You'll want to know: inquilīnus, inquilīnī}
\textbf{inquilīnus, inquilīnī} — \textquotedblleft{}a tenant, a lodger\textquotedblright{}.

\begin{quote}
Inquilīnus dominō mercēdem dat.
\end{quote}

\begin{quote}
\emph{The tenant pays the landlord his rent.}
\end{quote}

Second declension, masculine: \emph{inquilīnus, inquilīnī m.} Someone living in a house that is not his own; most people in a big city rented.

\begin{cousinweb}
\emph{inquilīnus} is built on \emph{incolō} 'to live in', from \emph{in-} and \emph{colō} 'to dwell, to till' (Lewis \& Short, from an older \emph{incolīnus}). EtymDB lists it without a parent.
\end{cousinweb}

\begin{grammarlens}[title={What you've built}]
One more word for telling, explaining and describing at B1.
\end{grammarlens}

\subsection*{Your turn}
\begin{itemize}
\raggedright
  \item \emph{Say it:} \emph{inquilīnus, inquilīnī}
  \item \emph{Say it:} \emph{inquilīnus, inquilīnī}, once more
  \item \emph{Say it:} a sentence of your own with \emph{inquilīnus, inquilīnī}
  \item \emph{Recall:} say the Latin for to delay, then the Latin for to stop for a while, then say \emph{inquilīnus, inquilīnī} again
\end{itemize}

\begin{tcolorbox}[breakable,colback=teal!4,colframe=teal!35!black,title={Before you move on}]
What does inquilīnus, inquilīnī mean? (\textquotedblleft{}A tenant, a lodger\textquotedblright{}.) Say it once more.
\end{tcolorbox}
\section[oppidānus, oppidānī]{oppidānus, oppidānī — a townsman, an inhabitant of a town}

\label{lesson:LA-C456-oppidanus}



\begin{quote}
Before the new one: say the Latin for to stop for a while, then the Latin for a tenant.
\end{quote}

\subsection*{You'll want to know: oppidānus, oppidānī}
\textbf{oppidānus, oppidānī} — \textquotedblleft{}a townsman, an inhabitant of a town\textquotedblright{}.

\begin{quote}
Oppidānī in forō mercātōrēs exspectant.
\end{quote}

\begin{quote}
\emph{The townspeople wait for the merchants in the forum.}
\end{quote}

Second declension, masculine: \emph{oppidānus, oppidānī m.}, mostly plural. Also an adjective: of a country town.

\begin{cousinweb}
\emph{oppidānus} is built on \emph{oppidum} 'town' with the suffix \emph{-ānus} 'belonging to' (EtymDB; Lewis \& Short).
\end{cousinweb}

\begin{grammarlens}[title={What you've built}]
One more word for telling, explaining and describing at B1.
\end{grammarlens}

\subsection*{Your turn}
\begin{itemize}
\raggedright
  \item \emph{Say it:} \emph{oppidānus, oppidānī}
  \item \emph{Say it:} \emph{oppidānus, oppidānī}, once more
  \item \emph{Say it:} a sentence of your own with \emph{oppidānus, oppidānī}
  \item \emph{Recall:} say the Latin for to stop for a while, then the Latin for a tenant, then say \emph{oppidānus, oppidānī} again
\end{itemize}

\begin{tcolorbox}[breakable,colback=teal!4,colframe=teal!35!black,title={Before you move on}]
What does oppidānus, oppidānī mean? (\textquotedblleft{}A townsman, an inhabitant of a town\textquotedblright{}.) Say it once more.
\end{tcolorbox}
\section[incola, incolae]{incola, incolae — a resident, an inhabitant}

\label{lesson:LA-C456-incola}



\begin{quote}
Before the new one: say the Latin for a tenant, then the Latin for a townsman.
\end{quote}

\subsection*{You'll want to know: incola, incolae}
\textbf{incola, incolae} — \textquotedblleft{}a resident, an inhabitant\textquotedblright{}.

\begin{quote}
Incolae vīcī nōs benignē accipiunt.
\end{quote}

\begin{quote}
\emph{The people who live in the village receive us kindly.}
\end{quote}

First declension, masculine or feminine: \emph{incola, incolae}. Like \emph{nauta} and \emph{agricola}, a first-declension noun for a person who is often a man.

\begin{cousinweb}
\emph{incola} is built on \emph{incolō} 'to live in' (EtymDB; Lewis \& Short), which comes from \emph{in-} and \emph{colō} 'to dwell'.
\end{cousinweb}

\begin{grammarlens}[title={What you've built}]
One more word for telling, explaining and describing at B1.
\end{grammarlens}

\subsection*{Your turn}
\begin{itemize}
\raggedright
  \item \emph{Say it:} \emph{incola, incolae}
  \item \emph{Say it:} \emph{incola, incolae}, once more
  \item \emph{Say it:} a sentence of your own with \emph{incola, incolae}
  \item \emph{Recall:} say the Latin for a tenant, then the Latin for a townsman, then say \emph{incola, incolae} again
\end{itemize}

\begin{tcolorbox}[breakable,colback=teal!4,colframe=teal!35!black,title={Before you move on}]
What does incola, incolae mean? (\textquotedblleft{}A resident, an inhabitant\textquotedblright{}.) Say it once more.
\end{tcolorbox}
\section[incolō, incolere]{incolō, incolere — to live in, to inhabit}

\label{lesson:LA-C456-incolo}



\begin{quote}
Before the new one: say the Latin for a townsman, then the Latin for a resident.
\end{quote}

\subsection*{You'll want to know: incolō, incolere}
\textbf{incolō, incolere} — \textquotedblleft{}to live in, to inhabit\textquotedblright{}.

\begin{quote}
Multī hominēs hanc urbem incolunt.
\end{quote}

\begin{quote}
\emph{Many people live in this city.}
\end{quote}

\emph{incolō, incolere, incoluī}, with an accusative of the place lived in.

\begin{cousinweb}
\emph{incolō} is built from \emph{in-} 'in' and \emph{colō} 'to dwell, to till' (EtymDB; Lewis \& Short).
\end{cousinweb}

\begin{grammarlens}[title={What you've built}]
One more word for telling, explaining and describing at B1.
\end{grammarlens}

\subsection*{Your turn}
\begin{itemize}
\raggedright
  \item \emph{Say it:} \emph{incolō, incolere}
  \item \emph{Say it:} \emph{incolō, incolere}, once more
  \item \emph{Say it:} a sentence of your own with \emph{incolō, incolere}
  \item \emph{Recall:} say the Latin for a townsman, then the Latin for a resident, then say \emph{incolō, incolere} again
\end{itemize}

\begin{tcolorbox}[breakable,colback=teal!4,colframe=teal!35!black,title={Before you move on}]
What does incolō, incolere mean? (\textquotedblleft{}To live in, to inhabit\textquotedblright{}.) Say it once more.
\end{tcolorbox}
\section[habitātor, habitātōris]{habitātor, habitātōris — an occupant, a dweller}

\label{lesson:LA-C456-habitator}



\begin{quote}
Before the new one: say the Latin for a resident, then the Latin for to live in.
\end{quote}

\subsection*{You'll want to know: habitātor, habitātōris}
\textbf{habitātor, habitātōris} — \textquotedblleft{}an occupant, a dweller\textquotedblright{}.

\begin{quote}
Domus vetus est, sed habitātōrēs eam amant.
\end{quote}

\begin{quote}
\emph{The house is old, but the people who live in it love it.}
\end{quote}

Third declension, masculine: \emph{habitātor, habitātōris m.}, from \emph{habitō} 'to live'. English \emph{habitat} comes from the same verb.

\begin{cousinweb}
\emph{habitātor} is built on \emph{habitō} 'to live, to dwell' with the suffix \emph{-tor} 'one who' (Lewis \& Short). \emph{habitō} itself is a frequentative of \emph{habeō} 'to have, to hold'.
\end{cousinweb}

\begin{grammarlens}[title={What you've built}]
One more word for telling, explaining and describing at B1.
\end{grammarlens}

\subsection*{Your turn}
\begin{itemize}
\raggedright
  \item \emph{Say it:} \emph{habitātor, habitātōris}
  \item \emph{Say it:} \emph{habitātor, habitātōris}, once more
  \item \emph{Say it:} a sentence of your own with \emph{habitātor, habitātōris}
  \item \emph{Recall:} say the Latin for a resident, then the Latin for to live in, then say \emph{habitātor, habitātōris} again
\end{itemize}

\begin{tcolorbox}[breakable,colback=teal!4,colframe=teal!35!black,title={Before you move on}]
What does habitātor, habitātōris mean? (\textquotedblleft{}An occupant, a dweller\textquotedblright{}.) Say it once more.
\end{tcolorbox}
